%===============================================================================
% APPENDIX: MATHEMATICAL PROOFS AND THEORETICAL BACKGROUND
%===============================================================================
\section{Mathematical Proofs and Theoretical Background}
\label{app:proofs}
%===============================================================================

This appendix provides a detailed theoretical examination of the core challenges in causal representation learning that motivate our proposed architecture.
We begin by formally defining and proving the incompatibility of dense encoders via the Intervention Spillover Problem, and then discuss the role of prior knowledge in differentiating causal modeling strategies.

%-------------------------------------------------------------------------------
% Subsection: The Intervention Spillover Problem
%-------------------------------------------------------------------------------
\subsection{The Intervention Spillover Problem}
\label{app:proofs:spillover}
%-------------------------------------------------------------------------------

This section provides the formal theoretical background for the \emph{Intervention Spillover Problem}.
We formally define this problem, arguing that it is a direct consequence of the inherent properties of dense encoder architectures commonly used in \gls{crl}.
We then prove the fundamental incompatibility between these architectures and the requirements for causal identifiability.

\begin{definition}[The Intervention Spillover Problem]
\label{def:app:intervention_spillover_problem}
The \emph{Intervention Spillover Problem} is the phenomenon where a targeted single-variable intervention in a high-dimensional observation space is transformed into a dense, multi-node signal in the lower-dimensional latent space.
This issue arises as a direct consequence of using a practical, dense encoder architecture, which inevitably corrupts the intervention's sparsity and creates a fundamental incompatibility with \gls{crl} methods that require sparse, single-node latent interventions for identifiability.
\end{definition}

%------------------------------------------------------------------------------
% Subsubsection: Mathematical Preliminaries and Setup
%------------------------------------------------------------------------------
\subsubsection{Mathematical Preliminaries and Setup}
\label{app:proofs:spillover:setup}
%------------------------------------------------------------------------------

While the problem described is general, this analysis formalizes it within the context of a \emph{linear \gls{scm}} for clarity.
Let the gene expression space be $\realnumbers^\observeddim$ and the latent space be $\realnumbers^\latentdim$, with $\latentdim \ll \observeddim$.
Let $\practicalencoder: \realnumbers^\observeddim \to \realnumbers^\latentdim$ be the practical, differentiable encoder network we train.
It maps an observation vector $\variables$ to an \emph{approximated} latent vector $\latentvars = \practicalencoder(\variables)$.
The local behavior of this encoder is described by its Jacobian matrix $\mJ_\theta(\variables) \in \realnumbers^{\latentdim \times \observeddim}$.

We model the latent space with a linear \gls{scm}, where causal relationships are defined by an adjacency matrix $\graphmatrix$.
An intervention on a single gene creates a post-intervention vector $\variables_{\text{int}}$ from a corresponding control vector $\variables_{\text{obs}}$.
This experimental paradigm is characteristic of modern pooled CRISPR screens (e.g., Perturb-seq), where the genetic target of each intervention is known by design.

%------------------------------------------------------------------------------
% Subsubsection: Formal Assumptions
%------------------------------------------------------------------------------
\subsubsection{Formal Assumptions}
\label{app:proofs:spillover:assumptions}
%------------------------------------------------------------------------------

\begin{assumption}[Ideal Causal Encoder and Single-Node Intervention Target]
\label{assump:app:ideal_encoder}
An \emph{Ideal Causal Encoder}, $\idealencoder$, is an unknown oracle function that perfectly disentangles the effect of any intervention.
The key property is that the \emph{true} latent change, $\latentchangeideal = \idealencoder(\variables_{\text{int}}) - \idealencoder(\variables_{\text{obs}})$, corresponds to a single-node perturbation.
This ideal latent change is defined as $\latentchangeideal = \maskvector \cdot \interventionstrength$, where $\interventionstrength$ is the latent intervention strength and the target mask $\maskvector \in \{0, 1\}^\latentdim$ is \emph{one-hot}:
\begin{equation}
\sum_{i=1}^{\latentdim} m_i = 1
\label{eq:app:one-hot}
\end{equation}
\end{assumption}

In a \emph{closed-world experimental setting}, where the targeted gene is known \apriori{}, the behavior of the $\idealencoder$ must fulfill the single-node intervention requirement.
The following lemma formalizes this ideal mapping.

\begin{lemma}[Equivalence in Ideal Latent Space]
\label{lemma:app:ideal_equivalence}
For an Ideal Causal Encoder $\idealencoder$ operating in a closed-world setting, the latent representation of the post-intervention sample is equivalent to the latent representation of the control sample plus the ideal single-node latent perturbation.
Formally:
\begin{equation}
    \idealencoder(\variables_{\text{int}}) = \idealencoder(\variables_{\text{obs}}) + \maskvector \cdot \interventionstrength
\end{equation}
\end{lemma}

\begin{proof}
The proof follows directly from the definition of the ideal latent change in \cref{assump:app:ideal_encoder}.
We start with the identity $\latentchangeideal = \idealencoder(\variables_{\text{int}}) - \idealencoder(\variables_{\text{obs}})$.
Substituting the single-node perturbation gives $\idealencoder(\variables_{\text{int}}) - \idealencoder(\variables_{\text{obs}}) = \maskvector \cdot \interventionstrength$.
Rearranging the terms yields the desired result.
\end{proof}

\begin{assumption}[Practical Model Approximation and Dense Jacobian]
\label{assump:app:practical_model}
Our \emph{Practical Model} uses an encoder $\practicalencoder(\variables)$ and a perturbation inference network $\interventionencoder$.
The underlying encoder network has a \emph{dense Jacobian matrix} $\mJ_\theta(\variables)$.
\end{assumption}

\begin{remark}[Architectural Bias vs. Learned Function]
It is important to distinguish between the architecture of the encoder and the function it learns.
While it is theoretically possible for a densely-connected network to learn an effectively sparse Jacobian, the architecture itself possesses a strong inductive bias towards dense mappings.
Standard training objectives provide no direct incentive for learning a sparse function.
\end{remark}

%------------------------------------------------------------------------------
% Subsubsection: Main Result: The Incompatibility
%------------------------------------------------------------------------------
\subsubsection{Main Result: The Incompatibility}
\label{app:proofs:spillover:lemma_proof}
%------------------------------------------------------------------------------

\begin{lemma}[Incompatibility of the Practical Model]
\label{lemma:app:incompatibility}
Given the practical model's reliance on a dense encoder network (\cref{assump:app:practical_model}), there is no guarantee that an intervention on a single gene will produce a single-node latent intervention target vector $\maskvector$ as required by the ideal model in \cref{assump:app:ideal_encoder}.
\end{lemma}
\begin{proof}
The proof rests on the direct conflict between the output of the practical encoder and the requirements of the ideal model in a closed-world setting.
In this setting, we have access to both the control sample $\variables_{\text{obs}}$ and the post-intervention sample $\variables_{\text{int}}$.
Our trained practical encoder, $\practicalencoder$, is a known deterministic function (e.g., a trained MLP).

When we process these known samples, the encoder produces an approximated latent change:
\begin{equation}
    \latentchangeapprox = \practicalencoder(\variables_{\text{int}}) - \practicalencoder(\variables_{\text{obs}})
\end{equation}
Due to the dense Jacobian of $\practicalencoder$ (\cref{assump:app:practical_model}), each component of the latent vector $\latentvars$ is a function of all input components in $\variables$.
Consequently, a change in a single gene will result in a change across nearly all latent dimensions.
The resulting vector $\latentchangeapprox$ is therefore dense.

The objective for a conditional network $\interventionencoder$ is to learn to predict this latent change.
For the model to perfectly capture the intervention effect, the learned perturbation $(\maskvector \cdot \interventionstrength)$ must match the observed latent change.
This implies minimizing the loss:
\begin{equation}
    \mathcal{L} = || \latentchangeapprox - (\maskvector \cdot \interventionstrength) ||^2
\end{equation}
To achieve a loss of zero, the learned perturbation $(\maskvector \cdot \interventionstrength)$ must be equal to the dense vector $\latentchangeapprox$.
This forces the learned target mask $\maskvector$ to be dense, which directly contradicts the one-hot requirement for a single-node intervention as defined in \cref{eq:app:one-hot}.
Therefore, the practical model's architecture is fundamentally misaligned with the goal of identifying single-node latent interventions.
\end{proof}

%------------------------------------------------------------------------------
% Subsubsection: Implications for Causal Discovery in Dense Encoders
%------------------------------------------------------------------------------
\subsubsection{Implications for Causal Discovery in Dense Encoders}
\label{app:proofs:spillover:implications}
%------------------------------------------------------------------------------

This mathematical reality creates a cascade of problems for causal discovery models that rely on a learned intervention network to identify the latent target of a perturbation.
The core of the issue is the need to make a discrete, single-target selection.
This is typically addressed using a temperature-scaled softmax function, a mechanism that mimics the Gumbel-Softmax trick for differentiable categorical sampling~\citep{jang2016categorical}.
Given the output logits, $\vo$, from an intervention network, the probability for each target is calculated as:
\begin{equation}
\pi_i = \frac{\exp(\vo_i / \gumbeltemp)}{\sum_{k=1}^{\latentdim} \exp(\vo_k / \gumbeltemp)}
\label{eq:app:temp_softmax}
\end{equation}
where $\gumbeltemp > 0$ is a temperature hyperparameter that controls the sharpness of the output distribution.

As the temperature $\gumbeltemp \to 0$, the softmax output approaches a discrete, one-hot vector (a proxy for a Dirac delta function), representing a single, "hard" choice.
However, this creates a fundamental contradiction when applied to the output of a dense encoder.
The encoder produces a dense latent change vector, $\latentchangeapprox$, where the evidence suggests a multi-node perturbation.
To satisfy the model's objective of selecting a single target, a low temperature $\gumbeltemp$ must be used to force the softmax into a "winner-take-all" state.
This forces the model to ignore the dense evidence and declare a single winner, which can lead to it learning spurious, causally unfaithful mappings.

The temperature-scaled softmax is therefore caught in an architectural dissonance: it is a mechanism forced to make a discrete choice that fundamentally misrepresents the dense, continuous signal it receives from the encoder.

Furthermore, in the closed-world setting of a Perturb-seq experiment, the mapping from a perturbed gene to its corresponding latent variable is known beforehand.
Therefore, the task of ``exploratory perturbation'' (i.e., discovering which latent variable was the target) is not actually necessary.
The true challenge is to learn the \emph{scale} or \emph{power} of the intervention's effect on the causal representation, not its target, a task for which these architectures are ill-suited.

%-------------------------------------------------------------------------------
% Subsection: The Role of Prior Knowledge in Causal Modeling
%-------------------------------------------------------------------------------
\subsection{The Role of Prior Knowledge in Causal Modeling}
\label{app:proofs:prior_knowledge}
%-------------------------------------------------------------------------------

The theoretical gaps in existing models can be understood by positioning them on a spectrum between two complementary modes of scientific inquiry: exploratory and confirmatory research~\citep{pmlr-v235-herrmann24b}.
Exploratory research is concerned with hypothesis generation; it involves analyzing data to uncover novel patterns, structures, and relationships without a strong pre-existing theory.
In contrast, confirmatory research is concerned with hypothesis testing; it begins with a specific, pre-formulated hypothesis and uses data to determine if that hypothesis is supported or refuted.
This distinction helps clarify the intended purpose and inherent trade-offs of different causal models.
For instance, \textbf{\gls{dvae}} is primarily exploratory, aiming to discover latent variables and their causal graph from scratch~\citep{zhang2023identifiability}.
In contrast, \textbf{SENA} takes a more targeted approach, acting as a confirmatory model for pathways while exploring perturbation effects; it \textbf{confirms} latent variables as known biological pathways and \textbf{explores} which pathway is targeted by a perturbation~\citep{de2025interpretable}.

The approach by \citet{de2025interpretable} highlights the close relationship between \gls{scm} and \gls{crl}.
By assuming its causal variables (the biological pathways) are known \emph{a priori}, it operates like a traditional \gls{scm}.
However, because it must still \emph{learn} a representation of how to map raw gene expression data onto these variables and then discover the causal graph between them, it firmly resides within the \gls{crl} paradigm.
It represents a specific, constrained form of \gls{crl} where prior knowledge heavily guides the representation learning process.
This confirmatory design choice is motivated by the goal of enhancing the biological interpretability of the learned causal model.

However, this reliance on predefined pathways introduces significant limitations, primarily stemming from the challenge of knowledge transfer.
Biological pathways exhibit strong cell-type and cell-line specificity, meaning that mechanisms characterized in one cellular context often do not generalize to another~\citep{schneider2017tissue,gamazon2018using,hekselman2020mechanisms}.
This context dependence poses a fundamental barrier to transferring pathway insights across different \emph{in vitro} models and ultimately to patient tissues~\citep{walter2019biological}.
Consequently, the findings from confirmatory models are difficult to apply to less-researched cell types, which limits the scalability and broader applicability of these approaches.
Furthermore, this confirmatory strategy may not be suitable for smaller-scale experiments that only incorporate a limited set of genes, making it difficult to map observations to comprehensive pathway definitions.

%-------------------------------------------------------------------------------
\subsection{Theoretical Necessity of Structured Pairing: Proof of Mean Collapse}
\label{app:theory:mean_collapse}
%-------------------------------------------------------------------------------

We formally justify the necessity of \glsentryshort{ot} pairing over the random pairing strategy common in baseline models (e.g., dVAE, SENA).

\paragraph{1. The Model's Goal}
Let $\variables \in \observsupport$ be a control cell and $\variables' \in \interventionaldist$ be a perturbed cell under intervention $p$. The model $\practicalencoder$ minimizes the reconstruction loss:
\begin{equation}
    \min_\theta \mathbb{E}_{\variables \sim \observsupport, \variables' \sim \interventionaldist} [ \| \variables' - \practicalencoder(\variables, p) \|^2 ]
\end{equation}

\paragraph{2. The Optimal Solution (MMSE)}
By the principles of Minimum Mean Squared Error (MMSE) estimation, the optimal function $\idealencoder(\variables, p)$ that minimizes this loss is the conditional expectation:
\begin{equation}
    \idealencoder(\variables, p) = \mathbb{E}[\variables' \mid \variables, P=p]
\end{equation}

\paragraph{3. The Case for Random Pairing}
When pairing is performed randomly, the target $\variables'$ is selected independently of the input $\variables$, given the perturbation label ($\variables' \perp \variables \mid P=p$). This independence eliminates the dependency of the expectation on the specific cell state $\variables$:
\begin{equation}
    \mathbb{E}[\variables' \mid \variables, P=p] = \mathbb{E}[\variables' \mid P=p] = \meanexpression_p
\end{equation}
where $\meanexpression_p$ is the population mean for perturbation $p$.

\paragraph{4. Conclusion (Mean Collapse)}
The optimal strategy for the model weights is to predict a constant output $\meanexpression_p$, effectively ignoring the biological variation in the input $\variables$. This sets the Jacobian $\nabla_\variables \predictionfunc \to 0$, precluding the discovery of heterogeneous (cell-specific) treatment effects. \glsentryshort{raptorgraph}'s \glsentryshort{ot} pairing ensures $I(\variables, \variables') > 0$, allowing the model to learn how specific initial states respond to interventions.

